\section{Benchmark Setup}

EvoDrive Lab uses procedurally generated top-down driving tracks parameterized by a track seed and a small number of geometry controls. Each episode is defined by a centerline, left and right walls, a start gate, a finish gate, and intermediate checkpoint gates. The benchmark is designed to remain lightweight enough for laptop-scale experimentation while still exposing a meaningful train-versus-test generalization problem.

Each car observes the environment through forward-facing ray sensors and motion-state features. The default setting uses five rays in the main comparison, and the sensor ablation increases this to seven rays. Policies output steering and throttle actions. Episodes terminate on crash or timeout. The current benchmark uses fixed train, validation, and test suites, allowing the same policies to be evaluated both on optimization tracks and on unseen held-out tracks.

The benchmark reports mean reward, mean completion, crash rate, and mean episode length. These values are currently the most stable and directly supported evaluation metrics in the codebase, and they form the basis of the quantitative analysis. In addition to the primary metrics, the reporting pipeline derives a generalization-gap view by comparing validation performance with held-out test performance for the same trained policy.
